% from ega1-3-fr.tex
\oldpage[I]{119}
fermé $y=\mathfrak{J}$ et le point générique $(0)$, l'ensemble $U$ réduit au
point générique étant donc un ouvert $U=\operatorname{Spec}(K)$ dans $Y$. Si
$X$ est un $A$-préschéma (autrement dit un $Y$-préschéma),
$X\otimes_A K=X'$ n'est autre que le préschéma induit par $X$ sur
$\psi^{-1}(U)$, en désignant par $\psi$ le morphisme structural $X\to Y$. En
particulier, si $\varphi$ est le morphisme structural $P\to Y$, un
sous-préschéma fermé $X'$ de $P'=\varphi^{-1}(U)$ est donc un sous-préschéma
(localement fermé) de $P$. Si $P$ est noethérien (par exemple si $A$ est
noethérien et $P$ de type fini sur $A$), il existe un plus petit
sous-préschéma fermé $X=\overline{X'}$ de $P$ qui majore $X'$ (9.5.10), et
$X'$ est le préschéma induit par $X$ sur l'ouvert
$\varphi^{-1}(U)\cap X$, donc est isomorphe à $X\otimes_A K$ (9.5.10).
\emph{L'immersion de $X'$ dans $P'=P\otimes_A K$ permet donc de façon
canonique de considérer $X'$ comme étant de la forme $X'=X\otimes_A K$, où
$X$ est un $A$-préschéma.} On peut alors considérer le préschéma réduit mod.
$\mathfrak{J}$, $X_0=X\otimes_A k$, qui n'est autre d'ailleurs que la fibre
$\psi^{-1}(y)$ du point fermé $y$. Jusqu'à présent, faute d'une terminologie
adéquate, on avait évité d'introduire explicitement le $A$-préschéma $X$. Il
convient cependant de noter que toutes les assertions faites habituellement
sur le préschéma «~réduit mod. $\mathfrak{J}$~» $X_0$ doivent être regardées
comme conséquences d'assertions plus complètes concernant $X$ lui-même, et ne
peuvent être formulées et comprises de façon satisfaisante qu'en les
interprétant ainsi. Il semble d'ailleurs que les hypothèses faites reviennent
toujours à des hypothèses sur $X$ lui-même (indépendamment de la donnée
préalable d'une immersion de $X'$ dans un $\mathbf{P}_K^r$), ce qui permet de
donner des énoncés plus intrinsèques.
\end{env}

\begin{env}[3.7.3]
\label{I.3.7.3-fr}
Signalons enfin un fait très particulier, qui a sans doute contribué à
retarder la clarification conceptuelle de la situation envisagée ici~: si
$A$ est un anneau de valuation discrète et si $X$ est \emph{propre sur $A$}
(ce qui est en fait le cas si $X$ est un sous-préschéma fermé d'un
$\mathbf{P}_A^r$, cf. II, 5.5.4), les points de $X$ à valeurs dans $A$ et les
points de $X'$ à valeurs dans $K$ sont en correspondance biunivoque
(II, 7.3.8). C'est pourquoi on a souvent cru démontrer des résultats
concernant $X'$, alors qu'en réalité on démontrait des énoncés concernant
$X$, et qui restent valables (sous cette forme) lorsqu'on ne suppose plus
l'anneau local de base de dimension 1.
\end{env}
